\chapter{概率论与贝叶斯推理}
\label{ch:prob}

% ════════════════════════════════════════════════════════════════
%  P0 数学基础部：概率论与贝叶斯推理。
%  为 P1 五章(状态估计/非线性优化/卡尔曼-ESKF/…)与后续 SLAM/规划/控制各部
%  打概率地基：概率空间→随机变量→期望/方差/协方差→独立与条件→贝叶斯定理→
%  高斯分布与多元高斯条件分布(滤波核心)→大数律/中心极限→KL 散度与信息量。
%  自包含独立记录：标准概率论(低幻觉)，每条定义/定理给条件与证明骨架，全写入正文。
%  教学层遵循 v5.0 理论教学规范（动机→反面→历史→理论→陷阱→练习 + 符号表/速查/排查）。
% ════════════════════════════════════════════════════════════════

\begin{practice}[前置自测]
开始之前，先试答下列问题。若已能流畅作答，可只速读本章表格与速查卡；若卡住，正是本章要为你解决的。
\begin{enumerate}
  \item 机器人测得一段含噪声的距离序列，你说“真值大概在 $3.0\pm0.1$ 米”。这句“大概”“$\pm$”在数学上究竟指什么对象？
  \item 已知传感器在物体存在时有 $99\%$ 概率报警（真阳率），但在没有物体时也有 $5\%$ 概率误报。某次报警了，物体真的存在的概率是多少？为什么不能直接说 $99\%$？
  \item 两个随机变量“不相关”和“相互独立”是一回事吗？给一个不相关但不独立的例子。
  \item 多元高斯 $\mathcal{N}(\boldsymbol{\mu},\boldsymbol{\Sigma})$ 的密度里出现 $(\mathbf{x}-\boldsymbol{\mu})^\top\boldsymbol{\Sigma}^{-1}(\mathbf{x}-\boldsymbol{\mu})$，它的等高面是什么形状？$\boldsymbol{\Sigma}^{-1}$ 在其中扮演什么角色？
  \item 给定联合高斯 $p(\mathbf{x},\mathbf{y})$，把 $\mathbf{y}$ 观测到具体值后，$\mathbf{x}$ 的条件分布 $p(\mathbf{x}\mid\mathbf{y})$ 还是高斯吗？它的均值与协方差怎么算？这条公式和卡尔曼滤波有什么关系？
  \item 把 $1000$ 个独立同分布、但\emph{并非}高斯的噪声样本求平均，结果的分布趋向什么？这为什么让我们敢在工程里到处假设高斯？
\end{enumerate}
\end{practice}

\section*{本章目标}
学完本章，你应当能够：
\begin{enumerate}
  \item \textbf{说清}“概率”是测度而非频率直觉，并写出概率空间 $(\Omega,\mathcal{F},\mathbb{P})$ 三要素，理解为何需要 $\sigma$-代数；
  \item \textbf{掌握}随机变量、分布函数、密度/质量函数，以及随机向量的联合/边缘/条件分布；
  \item \textbf{运用}期望的线性性、方差/协方差/协方差矩阵的定义与性质，独立证明线性变换 $\mathbf{y}=\mathbf{A}\mathbf{x}+\mathbf{b}$ 下均值与协方差的传播律；
  \item \textbf{区分}相互独立与不相关，举出二者不等价的反例，并说明高斯下二者何时合流；
  \item \textbf{独立推导}贝叶斯定理，理解先验/似然/后验/证据四者的角色，并用它解释“报警!=存在”的反直觉；
  \item \textbf{掌握}一元与多元高斯分布的密度、马氏距离的几何含义、高斯在线性变换与边缘化下的封闭性；
  \item \textbf{独立推导}多元高斯的\emph{条件分布}公式（滤波核心），并看清它就是卡尔曼增益的来源；
  \item \textbf{陈述并理解}大数定律与中心极限定理，解释它们为何是“蒙特卡洛”与“到处假设高斯”的理论依据；
  \item \textbf{运用} KL 散度、（微分）熵、互信息度量分布差异与信息量，理解“高斯下取负对数 $=$ 马氏二次型”这一贯穿全书估计章的桥；
  \item \textbf{定位}本章为后续——状态估计（\cref{ch:slam_est}）、卡尔曼与误差状态滤波（\cref{ch:eskf}）、SLAM 后端与因子图——提供统一的概率语言与全部前置工具。
\end{enumerate}

\subsection*{本章知识导航}
本章是一条“\emph{从“随机”这件事的严格定义，走到能为状态估计提供闭式工具}”的主线。结构如下：

\begin{center}
\small
\begin{tikzpicture}[
  node distance=4mm and 6mm, >=Stealth,
  blk/.style={draw, rounded corners, align=center, font=\small, inner sep=3pt, fill=main!6, draw=main!60!black},
  grp/.style={draw, rounded corners, align=center, font=\small, inner sep=3pt, fill=second!6, draw=second!60!black},
  ln/.style={->, thick, gray!60!black}]
\node[blk] (space) {概率空间\\ $(\Omega,\mathcal{F},\mathbb{P})$};
\node[blk, right=of space] (rv) {随机变量\\ 分布/密度};
\node[blk, right=of rv] (mom) {期望/方差\\ 协方差};
\node[grp, below=8mm of space] (cond) {独立与条件\\ 概率};
\node[grp, right=of cond] (bayes) {贝叶斯定理\\ 先验$\to$后验};
\node[grp, right=of bayes] (gauss) {高斯分布\\ 多元 $\mathcal{N}$};
\node[blk, below=8mm of cond, fill=third!8, draw=third!60!black] (condg) {高斯条件分布\\（滤波核心）};
\node[blk, right=of condg, fill=third!8, draw=third!60!black] (limit) {大数律 / CLT\\ KL 散度 / 熵};
\draw[ln] (space)--(rv); \draw[ln] (rv)--(mom);
\draw[ln] (mom.south) -- ++(0,-4.5mm) -| (cond.north);
\draw[ln] (cond)--(bayes); \draw[ln] (bayes)--(gauss);
\draw[ln] (gauss.south)-- (condg.north);
\draw[ln] (condg)--(limit);
\end{tikzpicture}
\end{center}

\noindent\textbf{推荐路径}：第 1--3 节（概率空间—随机变量—期望/协方差）是任何读者的主干，建立“随机量 $+$ 一阶/二阶矩”的最小工作集；第 4--5 节（独立/条件、贝叶斯定理）是估计与滤波的逻辑骨架——贝叶斯定理（\cref{thm:bayes}）是后续一切“由观测更新信念”的母定理；第 6 节高斯分布与第 7 节\emph{多元高斯条件分布}（\cref{thm:gauss-cond}）是全书最常用的闭式工具，卡尔曼滤波（\cref{ch:eskf}）的增益正是它的特例，强烈建议精读；第 8 节大数律/中心极限（\cref{thm:clt}）解释“为何到处敢用高斯、为何蒙特卡洛收敛”；第 9 节 KL 散度/熵/互信息（\cref{def:kl}）为信息形式估计、主动感知与变分推断提供尺子。带（选读）标记的测度论细节（$\sigma$-代数、可测性、勒贝格积分）只在需要严格性时回看，工程读者可先信其结论。

\subsection*{前置知识桥接}
本章只假定微积分（多元函数、重积分、泰勒展开）与线性代数（矩阵乘法、对称正定矩阵、特征分解、行列式、逆）——后者在 \cref{ch:linalg} 系统给出，其中\emph{对称正定矩阵的特征分解}是理解协方差椭球（\cref{sec:prob-gauss}）的关键，\emph{舒尔补与分块求逆}是高斯条件分布（\cref{sec:prob-gauss-cond}）的代数引擎。本章\emph{不}假定任何先前的概率知识，从概率的定义讲起。读者若已学过本科概率论，可把第 1--3 节当作“记号对齐 $+$ 把方差升级为协方差矩阵”的速读，重点放在第 5、7、9 节。

\subsection*{如果跳过本章会怎样}
\begin{itemize}
  \item 在 \cref{ch:slam_est}（状态估计）里，你会读到“SLAM 不是解方程，而是求后验 $p(\mathbf{x}\mid\mathbf{z})$ 的众数”，却\emph{不知道}后验、似然、先验各指什么，也无法理解“高斯噪声下最大化概率为何等于最小化误差平方和”；
  \item 在 \cref{ch:eskf}（卡尔曼 / ESKF）里，预测—更新两步会反复用到“两个高斯相乘还是高斯”“线性变换下高斯封闭”“条件高斯的均值线性插值”，这些若不先证清，卡尔曼增益 $\mathbf{K}$ 就只能当公式背、无法调试；
  \item 在因子图、主动感知、变分 SLAM 里会出现熵、互信息、KL 散度作为目标函数，不懂它们就只能照抄目标式而不知所求为何。
\end{itemize}

\begin{table}[H]\centering\small
\caption{本章阅读方式建议}
\begin{tabular}{lll}
\toprule
阅读方式 & 时间 & 适合谁 \\
\midrule
精读（含全部推导与练习） & 4--6 小时 & 首次系统学习、要做估计/SLAM 后端推导 \\
速读（跳过测度论细节与冗长证明） & 1.5 小时 & 有概率基础、想对齐本书记号与协方差工具 \\
速查（只看小结的符号表 + 定理速查表） & 15 分钟 & 写估计代码时回来查公式 \\
\bottomrule
\end{tabular}
\end{table}

\begin{note}
\textbf{本章记号与约定.}\;
随机变量用大写 $X,Y$，其取值（实现/样本）用小写 $x,y$；随机\emph{向量}沿用全书黑体并大写其名时仍写 $\mathbf{x}$（如 $\mathbf{x}\sim\mathcal{N}(\boldsymbol{\mu},\boldsymbol{\Sigma})$，与 \cref{ch:slam_est,ch:eskf} 一致，避免与状态记号冲突），上下文区分随机量与其实现。
概率测度记 $\mathbb{P}$，期望 $\mathbb{E}[\cdot]$，方差 $\operatorname{Var}[\cdot]$，协方差 $\operatorname{Cov}[\cdot,\cdot]$，协方差\emph{矩阵}记 $\boldsymbol{\Sigma}$、其逆（信息/精度矩阵）记 $\boldsymbol{\Lambda}=\boldsymbol{\Sigma}^{-1}$。
连续随机变量的密度记 $p(\cdot)$，离散的质量函数也记 $p(\cdot)$（按上下文区分），条件密度 $p(x\mid y)$。
高斯（正态）分布记 $\mathcal{N}(\boldsymbol{\mu},\boldsymbol{\Sigma})$（以\emph{协方差}而非精度参数化，与十四讲\cite{gaoxiang2019slam14}、Barfoot\cite{barfoot2024state} 一致）。
向量/矩阵粗体 $\mathbf{x},\mathbf{A}$，单位阵 $\mathbf{I}$；记号与 \cref{ch:lie}（李群）、\cref{ch:rigid_body}（刚体）保持一致（$\SO(3),\SE(3),\mathrm{Exp},\mathrm{Log}$）。
\end{note}

% ════════════════════════════════════════════════════════════════
\section{为什么机器人学需要概率论}
\label{sec:prob-why}

\paragraph{动机：机器人活在一个充满噪声的世界。}
机器人对世界的一切了解都来自\emph{传感器}，而没有任何传感器是精确的：激光雷达有测距噪声，相机有像素噪声与标定误差，IMU 有偏置与随机游走，轮速计打滑。与此同时，机器人对自身\emph{运动}的执行也不精确：电机有误差，地面有起伏。于是机器人面对的根本困境是：\textbf{它永远无法确切知道自己在哪、世界长什么样，只能基于含噪信息做出“最可能”的推断，并诚实地说出这个推断有多不确定}。把“不确定”这件事数学化、可计算化，正是概率论之于机器人学的意义。

\paragraph{反面：若假装一切确定会怎样。}
设想忽略噪声、把每次传感器读数都当作真值。轮速计积分给出里程，但每步的微小打滑误差会\emph{累积}：走一百米后位置可能偏差好几米，而你的程序却“自信满满”地认为自己分毫不差——这种\textbf{过度自信}比误差本身更危险，因为它让系统在错误的地方做出果断而错误的决策。反过来，若因噎废食、对一切都不敢相信，则信息无法融合，机器人寸步难行。概率论提供的正是这两极之间的\emph{定量平衡}：用协方差诚实地携带不确定度，用贝叶斯定理把多源含噪信息\emph{最优地}融合。

\paragraph{历史：从赌博到拉普拉斯到现代机器人。}
概率论起于 $17$ 世纪 Pascal 与 Fermat 关于赌博的通信；$18$ 世纪 Bayes 与 Laplace 把它发展为\emph{由观测更新信念}的推理工具——Laplace 用它估计天体轨道，正是“从含噪观测估计未知状态”的先声。$20$ 世纪 Kolmogorov 以测度论给出公理化基础\cite{bishop2006prml}，使概率成为严格的数学分支。机器人学自 $1990$ 年代起系统地以概率为统一语言：Thrun、Burgard、Fox 的《Probabilistic Robotics》\cite{thrun2005probabilistic} 把“一切皆后验”确立为范式，Barfoot\cite{barfoot2024state}、SLAM Handbook 与现代 VIO/SLAM 系统无不以贝叶斯滤波与因子图为内核。

\paragraph{核心思想。}
把机器人对世界的\emph{信念}表示为概率分布；新观测到来时，用\textbf{贝叶斯定理}把先验信念更新为后验信念；当一切（噪声、信念）都近似为\emph{高斯}时，这套更新有\emph{闭式解}——这就是卡尔曼滤波家族。本章自下而上地搭起这座大厦：先把“随机”严格定义（概率空间、随机变量），再装上度量不确定度的工具（期望、协方差），然后给出推理的母定理（贝叶斯），最后聚焦机器人学最常用的分布（高斯）及其能闭式运算的核心性质（线性封闭、条件分布）。

\begin{insight}[概率是“被压缩的信念”，不是“多次重复的频率”]
工程上常把概率直觉为“重复实验的频率”。这在能重复的场景（掷骰子）下没错，但机器人“此刻在哪”无法重复一万次。更有用的视角是\emph{贝叶斯}的：概率是对\textbf{一个命题为真的信念程度}的定量编码，可随证据更新。后续所有滤波/估计都站在这个视角上——它把“我有多确定”变成一个可加、可乘、可优化的数。两种视角在数学上由同一套公理（\cref{sec:prob-space}）支撑，不必纠结哲学，只需会算。
\end{insight}

% ════════════════════════════════════════════════════════════════
\section{概率空间与 \texorpdfstring{$\sigma$}{sigma}-代数}
\label{sec:prob-space}

\paragraph{动机：什么是“一个事件的概率”？}
要谈“概率”，先要说清\emph{什么东西}有概率。直觉上，概率是赋予“事件”的一个 $[0,1]$ 数。Kolmogorov 的洞见是：把全部可能结果、可谈论的事件、以及赋概率的规则三者一并形式化，概率论的一切便有了坚实地基。

\begin{definition}[概率空间]\label{def:prob-space}
概率空间是三元组 $(\Omega,\mathcal{F},\mathbb{P})$：
\begin{enumerate}
  \item \textbf{样本空间} $\Omega$：一次随机实验所有可能\emph{结果} $\omega$ 的集合；
  \item \textbf{事件域} $\mathcal{F}$：$\Omega$ 的某些子集构成的集族，每个 $A\in\mathcal{F}$ 称为\emph{事件}，要求 $\mathcal{F}$ 是 $\sigma$-代数（\cref{def:sigma-algebra}）；
  \item \textbf{概率测度} $\mathbb{P}:\mathcal{F}\to[0,1]$，满足 Kolmogorov 三公理：
  \begin{itemize}
    \item 非负：$\mathbb{P}(A)\ge0$；
    \item 规范：$\mathbb{P}(\Omega)=1$；
    \item 可数可加：若 $\{A_i\}_{i\ge1}$ 两两不交，则 $\mathbb{P}\big(\bigcup_i A_i\big)=\sum_i\mathbb{P}(A_i)$。
  \end{itemize}
\end{enumerate}
\end{definition}

\begin{definition}[$\sigma$-代数]\label{def:sigma-algebra}
$\Omega$ 上的集族 $\mathcal{F}\subseteq2^{\Omega}$ 是 $\sigma$-代数，若：(i) $\Omega\in\mathcal{F}$；(ii) 对补封闭：$A\in\mathcal{F}\Rightarrow A^{\mathrm c}\in\mathcal{F}$；(iii) 对\emph{可数}并封闭：$A_i\in\mathcal{F}\Rightarrow\bigcup_i A_i\in\mathcal{F}$。
\end{definition}

\paragraph{为什么需要 $\sigma$-代数？（选读）}
当 $\Omega$ 有限或可数时，可直接取 $\mathcal{F}=2^{\Omega}$（全部子集），$\sigma$-代数无非走过场。但当 $\Omega$ 不可数（如 $\mathbb{R}$，机器人位置取连续值）时，会出现\emph{无法一致赋概率}的“病态”子集（Vitali 集即一例），强行对\emph{所有}子集定义满足三公理的 $\mathbb{P}$ 会导致矛盾。出路是：只对一个“足够大但不致病态”的子集族 $\mathcal{F}$ 赋概率——这正是 $\sigma$-代数。$\mathbb{R}$ 上的标准选择是\emph{博雷尔} $\sigma$-代数 $\mathcal{B}(\mathbb{R})$，由全体开区间生成，足以涵盖一切工程上关心的事件（“位置落在 $[a,b]$”）。三条封闭性保证“事件的补、可数并、可数交、差”仍是事件，使概率运算（如 $\mathbb{P}(A^{\mathrm c})=1-\mathbb{P}(A)$）合法。

\begin{note}
\textbf{工程读者怎么用这一节.}\;
本节的全部目的，是让“连续随机变量的概率”在数学上站得住脚。实际计算中你几乎从不直接操作 $\sigma$-代数：连续量用\emph{密度} $p(x)$、离散量用\emph{质量函数}，事件概率由积分/求和给出（\cref{sec:prob-rv}）。可以把 $(\Omega,\mathcal{F},\mathbb{P})$ 当作背后的“记账系统”——保证一切自洽，但日常账目用密度来记。Barfoot\cite{barfoot2024state} 与十四讲\cite{gaoxiang2019slam14} 在正文中均默认这一层、直接从密度起步，本书亦然，仅在此交代地基。
\end{note}

\begin{proposition}[概率的基本运算法则]\label{prop:prob-rules}
由 \cref{def:prob-space} 可推出（设 $A,B\in\mathcal{F}$）：
\begin{enumerate}
  \item $\mathbb{P}(\varnothing)=0$，$\mathbb{P}(A^{\mathrm c})=1-\mathbb{P}(A)$；
  \item 单调性：$A\subseteq B\Rightarrow\mathbb{P}(A)\le\mathbb{P}(B)$；
  \item 容斥（加法公式）：$\mathbb{P}(A\cup B)=\mathbb{P}(A)+\mathbb{P}(B)-\mathbb{P}(A\cap B)$。
\end{enumerate}
\end{proposition}
\begin{proof}
由 $\Omega=A\cup A^{\mathrm c}$ 不交且 $\mathbb{P}(\Omega)=1$ 得 $\mathbb{P}(A^{\mathrm c})=1-\mathbb{P}(A)$；取 $A=\Omega$ 得 $\mathbb{P}(\varnothing)=0$。单调性：$B=A\cup(B\setminus A)$ 不交，故 $\mathbb{P}(B)=\mathbb{P}(A)+\mathbb{P}(B\setminus A)\ge\mathbb{P}(A)$。容斥：把 $A\cup B$ 拆为三块不交集 $A\setminus B$、$A\cap B$、$B\setminus A$，分别加可加性并配项即得。
\end{proof}

\begin{pitfall}[概念误区：“概率就是频率”，所以单次事件无概率]
“机器人现在恰好在 $(3,5)$ 处”不可重复，于是有人断言它“没有概率”。这是把概率\emph{狭义}地等同于长期频率。\cref{def:prob-space} 的公理\emph{不}规定 $\mathbb{P}$ 的哲学解释——它既可解释为频率，也可解释为\emph{信念程度}（贝叶斯派）。机器人学普遍采用后者：$\mathbb{P}(\text{在 }(3,5))$ 表示“基于当前所有证据，我相信它在那里的程度”。两种解释共享同一套运算法则，因此本书不纠缠哲学，只要求你会用公理算。
\end{pitfall}

\begin{practice}
\begin{enumerate}
  \item 验证 \cref{prop:prob-rules} 的容斥公式可推广到三事件：$\mathbb{P}(A\cup B\cup C)=\sum\mathbb{P}(\cdot)-\sum\mathbb{P}(\cdot\cap\cdot)+\mathbb{P}(A\cap B\cap C)$。
  \item （概念）在 $\Omega=\{1,2,3,4,5,6\}$（掷骰）上写出由事件“偶数”生成的最小 $\sigma$-代数，列出它的全部元素。
  \item （选读）解释为何当 $\Omega=\mathbb{R}$ 时不能取 $\mathcal{F}=2^{\mathbb{R}}$ 并对所有子集赋满足三公理的概率——查阅 Vitali 集构造的要点。
\end{enumerate}
\end{practice}

% ════════════════════════════════════════════════════════════════
\section{随机变量与分布}
\label{sec:prob-rv}

\paragraph{动机：把“结果”变成“数”。}
机器人关心的往往不是抽象结果 $\omega$，而是一个\emph{数值}：测得的距离、位姿的某分量、特征点的像素坐标。随机变量正是把样本空间映到实数（或实向量）的函数，使我们能用微积分处理随机性。

\begin{definition}[随机变量与随机向量]\label{def:random-variable}
$(\Omega,\mathcal{F},\mathbb{P})$ 上的\textbf{随机变量}是函数 $X:\Omega\to\mathbb{R}$，且对任意 $a\in\mathbb{R}$，集合 $\{\omega:X(\omega)\le a\}\in\mathcal{F}$（\emph{可测性}：保证“$X\le a$”是个有概率的事件）。$n$ 个随机变量叠成的 $\mathbf{x}=(X_1,\dots,X_n)^\top:\Omega\to\mathbb{R}^n$ 称为\textbf{随机向量}。
\end{definition}

\begin{definition}[分布函数、密度与质量函数]\label{def:cdf-pdf}
随机变量 $X$ 的\textbf{累积分布函数}（CDF）为 $F_X(x)=\mathbb{P}(X\le x)$，它单调不减、右连续、$F(-\infty)=0$、$F(+\infty)=1$。
\begin{itemize}
  \item 若 $X$ \emph{离散}（取值可数 $\{x_k\}$），其\textbf{概率质量函数}（PMF）为 $p(x_k)=\mathbb{P}(X=x_k)$，满足 $p(x_k)\ge0$、$\sum_k p(x_k)=1$；
  \item 若 $X$ \emph{连续}且存在非负可积函数 $p(x)$ 使 $F_X(x)=\int_{-\infty}^{x}p(t)\,dt$，则 $p$ 为\textbf{概率密度函数}（PDF），满足 $p(x)\ge0$、$\int_{\mathbb{R}}p(x)\,dx=1$，且 $\mathbb{P}(X\in[a,b])=\int_a^b p(x)\,dx$。
\end{itemize}
\end{definition}

\begin{pitfall}[概念误区：把连续分布的密度值 $p(x)$ 当成概率]
对连续 $X$，$p(x)$ \emph{不是}概率，而是\textbf{概率密度}——单点概率 $\mathbb{P}(X=x)=0$，只有落入\emph{区间}才有非零概率 $\int_a^b p$。因此 $p(x)$ 可以\emph{大于 $1$}（例如窄高斯的峰值），这毫不矛盾：密度乘以宽度才是概率。记住量纲：一维密度的量纲是“$1/$自变量单位”。把 $p(x)$ 误当概率，会在归一化、贝叶斯更新里酿成系统性错误。
\end{pitfall}

\paragraph{随机向量的联合、边缘与条件分布。}
机器人状态通常是\emph{向量}（位姿、速度、地图），故须处理多个随机变量的\emph{联合}行为。设随机向量 $(\mathbf{x},\mathbf{y})$ 有联合密度 $p(\mathbf{x},\mathbf{y})$，则：
\begin{align}
  \text{边缘}:\quad & p(\mathbf{x})=\int p(\mathbf{x},\mathbf{y})\,d\mathbf{y},\label{eq:marginal}\\
  \text{条件}:\quad & p(\mathbf{x}\mid\mathbf{y})=\frac{p(\mathbf{x},\mathbf{y})}{p(\mathbf{y})}\quad(p(\mathbf{y})>0).\label{eq:conditional}
\end{align}
\emph{边缘化}（\cref{eq:marginal}）是“把不关心的变量积分掉”——例如从“位姿与地图的联合”里只取位姿；\emph{条件化}（\cref{eq:conditional}）是“把已观测量固定、看剩余量的分布”——例如把观测 $\mathbf{y}$ 固定后看状态 $\mathbf{x}$ 的分布。这两步是状态估计的核心操作，\cref{ch:slam_est} 会反复用到；在高斯情形它们都有闭式（\cref{sec:prob-gauss-cond}）。

\begin{definition}[全概率公式]\label{def:total-prob}
若 $\{B_i\}$ 是 $\Omega$ 的一个\emph{划分}（两两不交、并为 $\Omega$，$\mathbb{P}(B_i)>0$），则对任意事件 $A$，
\begin{equation}
  \mathbb{P}(A)=\sum_i\mathbb{P}(A\mid B_i)\,\mathbb{P}(B_i),
  \qquad\text{（连续版）}\quad p(\mathbf{x})=\int p(\mathbf{x}\mid\mathbf{y})\,p(\mathbf{y})\,d\mathbf{y}.
  \label{eq:total-prob}
\end{equation}
\end{definition}
\noindent 全概率公式（\cref{eq:total-prob}）正是 \cref{eq:marginal} 与 \cref{eq:conditional} 的合体：先按 $\mathbf{y}$ 的取值“分情形”给出 $p(\mathbf{x}\mid\mathbf{y})$，再以 $p(\mathbf{y})$ 为权\emph{加权平均}回 $p(\mathbf{x})$。它是贝叶斯滤波\emph{预测步}（Chapman--Kolmogorov）的母公式：把上一时刻状态分布经运动模型“传播”到当前，正是一次全概率积分（详见 \cref{ch:eskf}）。

\begin{example}[一维到二维的直觉]\label{ex:joint-marginal}
设 $X$ 为机器人前向位置、$Y$ 为侧向位置，联合密度集中在一条斜带上（即 $X,Y$ 相关）。把 $Y$ 积掉得边缘 $p(x)$，相当于把二维“云”投影到 $x$ 轴；而若已测得侧向 $Y=y_0$，则 $p(x\mid y_0)$ 是把这条云沿 $y=y_0$ 切一刀、再归一化得到的细长分布——\emph{比边缘更窄}，因为多了一条信息。这正是“观测降低不确定度”的几何写照。
\end{example}

\begin{practice}
\begin{enumerate}
  \item 验证：由 \cref{def:cdf-pdf}，连续 $X$ 有 $p(x)=\dfrac{dF_X(x)}{dx}$（在 $F$ 可导处）。
  \item 设 $(\mathbf{x},\mathbf{y})$ 联合密度 $p(\mathbf{x},\mathbf{y})$ 已知，写出由 \cref{eq:marginal,eq:conditional} 得 $p(\mathbf{y}\mid\mathbf{x})$ 的步骤，并说明它与 $p(\mathbf{x}\mid\mathbf{y})$ 一般\emph{不}相等。
  \item （概念）举例说明密度 $p(x)>1$ 不违反任何公理（提示：取均匀分布 $U[0,0.5]$）。
\end{enumerate}
\end{practice}

% ════════════════════════════════════════════════════════════════
\section{期望、方差与协方差}
\label{sec:prob-moments}

\paragraph{动机：用几个数概括一整个分布。}
完整分布 $p(\mathbf{x})$ 信息丰富但难直接操作。工程上常用少数\emph{矩}来概括它：一阶矩（期望/均值）给出“中心在哪”，二阶中心矩（方差/协方差）给出“散布多大、各分量怎样联动”。对最常用的高斯分布，均值与协方差\emph{完全}决定整个分布——这正是滤波只需传播这两个量的根本原因。

\begin{definition}[期望]\label{def:expectation}
随机变量 $X$ 的\textbf{期望}（均值）为
\begin{equation}
  \mathbb{E}[X]=\begin{cases}\sum_k x_k\,p(x_k),&\text{离散};\\[2pt]\displaystyle\int_{\mathbb{R}}x\,p(x)\,dx,&\text{连续}.\end{cases}
  \label{eq:expectation}
\end{equation}
对随机向量 $\mathbf{x}$，$\mathbb{E}[\mathbf{x}]$ 逐分量取期望，得均值向量 $\boldsymbol{\mu}=\mathbb{E}[\mathbf{x}]$。函数 $g(X)$ 的期望为 $\mathbb{E}[g(X)]=\int g(x)\,p(x)\,dx$（\emph{无意识统计学家律}，LOTUS：无需先求 $g(X)$ 的分布）。
\end{definition}

\begin{theorem}[期望的线性性]\label{thm:linearity}
对任意随机变量 $X,Y$、常数 $a,b$，无论是否独立，
\begin{equation}
  \mathbb{E}[aX+bY]=a\,\mathbb{E}[X]+b\,\mathbb{E}[Y];\qquad
  \mathbb{E}[\mathbf{A}\mathbf{x}+\mathbf{b}]=\mathbf{A}\,\mathbb{E}[\mathbf{x}]+\mathbf{b}.
  \label{eq:linearity}
\end{equation}
\end{theorem}
\begin{proof}
由 \cref{eq:expectation} 与积分（求和）的线性性直接得标量式；向量式逐分量套标量式即可。关键在于线性性\emph{不要求独立}——这使它成为最常用、最稳的工具。
\end{proof}

\begin{definition}[方差、协方差与协方差矩阵]\label{def:covariance}
随机变量 $X$ 的\textbf{方差}与两变量的\textbf{协方差}：
\begin{align}
  \operatorname{Var}[X]&=\mathbb{E}\big[(X-\mathbb{E}[X])^2\big]=\mathbb{E}[X^2]-(\mathbb{E}[X])^2,\label{eq:variance}\\
  \operatorname{Cov}[X,Y]&=\mathbb{E}\big[(X-\mathbb{E}[X])(Y-\mathbb{E}[Y])\big]=\mathbb{E}[XY]-\mathbb{E}[X]\mathbb{E}[Y].\label{eq:cov}
\end{align}
随机向量 $\mathbf{x}$ 的\textbf{协方差矩阵}为
\begin{equation}
  \boldsymbol{\Sigma}=\operatorname{Cov}[\mathbf{x}]=\mathbb{E}\big[(\mathbf{x}-\boldsymbol{\mu})(\mathbf{x}-\boldsymbol{\mu})^\top\big]\in\mathbb{R}^{n\times n},
  \quad \Sigma_{ij}=\operatorname{Cov}[X_i,X_j].
  \label{eq:cov-matrix}
\end{equation}
对角元 $\Sigma_{ii}=\operatorname{Var}[X_i]$，非对角元刻画分量间的线性联动。
\end{definition}

\begin{proposition}[协方差矩阵对称半正定]\label{prop:cov-psd}
任意随机向量的协方差矩阵 $\boldsymbol{\Sigma}$ 对称且半正定（$\boldsymbol{\Sigma}=\boldsymbol{\Sigma}^\top\succeq0$）。
\end{proposition}
\begin{proof}
对称性显然（$\operatorname{Cov}[X_i,X_j]=\operatorname{Cov}[X_j,X_i]$）。半正定：对任意常向量 $\mathbf{a}$，标量 $s=\mathbf{a}^\top(\mathbf{x}-\boldsymbol{\mu})$ 满足
\[
  \mathbf{a}^\top\boldsymbol{\Sigma}\mathbf{a}=\mathbf{a}^\top\mathbb{E}[(\mathbf{x}-\boldsymbol{\mu})(\mathbf{x}-\boldsymbol{\mu})^\top]\mathbf{a}=\mathbb{E}[s^2]\ge0.
\]
若 $\boldsymbol{\Sigma}$ 还\emph{正定}（无方向上方差为零，即各分量非退化、不线性相关），则可逆——这是定义高斯密度（\cref{sec:prob-gauss}）所需。
\end{proof}

\begin{theorem}[线性变换的均值与协方差传播律]\label{thm:affine-prop}
设 $\mathbf{x}$ 有均值 $\boldsymbol{\mu}$、协方差 $\boldsymbol{\Sigma}$，$\mathbf{y}=\mathbf{A}\mathbf{x}+\mathbf{b}$（$\mathbf{A}$ 常矩阵、$\mathbf{b}$ 常向量），则
\begin{equation}
  \mathbb{E}[\mathbf{y}]=\mathbf{A}\boldsymbol{\mu}+\mathbf{b},\qquad
  \operatorname{Cov}[\mathbf{y}]=\mathbf{A}\boldsymbol{\Sigma}\mathbf{A}^\top.
  \label{eq:affine-prop}
\end{equation}
\end{theorem}
\begin{proof}
均值由线性性 \cref{eq:linearity} 立得。协方差：$\mathbf{y}-\mathbb{E}[\mathbf{y}]=\mathbf{A}(\mathbf{x}-\boldsymbol{\mu})$，故
\[
  \operatorname{Cov}[\mathbf{y}]=\mathbb{E}\big[\mathbf{A}(\mathbf{x}-\boldsymbol{\mu})(\mathbf{x}-\boldsymbol{\mu})^\top\mathbf{A}^\top\big]
  =\mathbf{A}\,\mathbb{E}[(\mathbf{x}-\boldsymbol{\mu})(\mathbf{x}-\boldsymbol{\mu})^\top]\,\mathbf{A}^\top=\mathbf{A}\boldsymbol{\Sigma}\mathbf{A}^\top,
\]
其中常矩阵 $\mathbf{A}$ 可提到期望外。注意平移 $\mathbf{b}$ \emph{不}影响协方差（协方差只看“相对中心的散布”）。
\end{proof}

\begin{insight}[$\mathbf{A}\boldsymbol{\Sigma}\mathbf{A}^\top$ 是全书最常用的“不确定度搬运”公式]
\cref{eq:affine-prop} 看似平凡，却是\emph{不确定度传播}的引擎：传感器模型把状态线性映为观测，$\mathbf{A}\boldsymbol{\Sigma}\mathbf{A}^\top$ 就把状态的协方差搬到观测空间；卡尔曼预测步把协方差经状态转移矩阵 $\mathbf{F}$ 传播为 $\mathbf{F}\boldsymbol{\Sigma}\mathbf{F}^\top$（\cref{ch:eskf}）；李群上把切空间协方差经伴随 $\mathrm{Ad}$ 搬运（\cref{ch:lie}）也是同一形式。非线性时，先一阶线性化（雅可比 $\mathbf{J}$ 充当 $\mathbf{A}$）再套此式，即得 EKF 的协方差传播 $\mathbf{J}\boldsymbol{\Sigma}\mathbf{J}^\top$。记住这一个式子，半部状态估计的“协方差怎么变”就有了答案。
\end{insight}

\paragraph{协方差的可加性。}对\emph{独立}（或仅\emph{不相关}）随机向量 $\mathbf{x},\mathbf{y}$，
\begin{equation}
  \operatorname{Cov}[\mathbf{x}+\mathbf{y}]=\operatorname{Cov}[\mathbf{x}]+\operatorname{Cov}[\mathbf{y}]\qquad(\mathbf{x},\mathbf{y}\text{ 不相关}).
  \label{eq:cov-additive}
\end{equation}
证明：展开 $\operatorname{Cov}[\mathbf{x}+\mathbf{y}]=\operatorname{Cov}[\mathbf{x}]+\operatorname{Cov}[\mathbf{y}]+\operatorname{Cov}[\mathbf{x},\mathbf{y}]+\operatorname{Cov}[\mathbf{y},\mathbf{x}]$，不相关时交叉项为零。这解释了为何独立噪声叠加时方差\emph{相加}（标准差按平方和开方），是误差预算（error budget）的基础。

\begin{practice}
\begin{enumerate}
  \item 由 \cref{eq:variance} 证明 $\operatorname{Var}[aX+b]=a^2\operatorname{Var}[X]$（平移不变、缩放平方）。
  \item 设 $\mathbf{x}\sim$ 协方差 $\boldsymbol{\Sigma}$，求 $\mathbf{y}=[X_1+X_2,\ X_1-X_2]^\top$ 的协方差矩阵，写出对应的 $\mathbf{A}$ 并套 \cref{eq:affine-prop}。
  \item 证明 \cref{prop:cov-psd}：若存在非零 $\mathbf{a}$ 使 $\mathbf{a}^\top\boldsymbol{\Sigma}\mathbf{a}=0$，则 $\boldsymbol{\Sigma}$ 奇异，且 $\mathbf{x}$ 落在一个低维仿射子空间上（退化分布）。
\end{enumerate}
\end{practice}

% ════════════════════════════════════════════════════════════════
\section{独立性、条件概率与贝叶斯定理}
\label{sec:prob-bayes}

\paragraph{动机：信息如何融合、信念如何更新。}
机器人不断获得新观测，需要把它们\emph{融合}进已有信念。这要求两件事：一是判断哪些信息源\emph{独立}（可分别处理、协方差相加），二是有一条\emph{由观测更新信念}的法则。后者就是贝叶斯定理——本章乃至全书估计部分的逻辑心脏。

\begin{definition}[条件概率与独立性]\label{def:independence}
事件 $A$ 在 $B$（$\mathbb{P}(B)>0$）下的\textbf{条件概率} $\mathbb{P}(A\mid B)=\dfrac{\mathbb{P}(A\cap B)}{\mathbb{P}(B)}$。事件 $A,B$ \textbf{独立}，若 $\mathbb{P}(A\cap B)=\mathbb{P}(A)\mathbb{P}(B)$，等价地 $\mathbb{P}(A\mid B)=\mathbb{P}(A)$（“知道 $B$ 不改变对 $A$ 的判断”）。随机变量 $X,Y$ 独立，若联合密度可分离 $p(x,y)=p(x)\,p(y)$（等价 $p(x\mid y)=p(x)$）。
\end{definition}

\begin{definition}[条件独立]\label{def:cond-independence}
$X,Y$ 在给定 $Z$ 下\textbf{条件独立}（记 $X\perp Y\mid Z$），若 $p(x,y\mid z)=p(x\mid z)\,p(y\mid z)$。这是马尔可夫性与因子图的代数核心：状态估计中“给定当前状态，过去与未来观测条件独立”正是此结构（\cref{ch:slam_est,ch:eskf}）。
\end{definition}

\begin{theorem}[贝叶斯定理]\label{thm:bayes}
对事件 $A,B$（$\mathbb{P}(B)>0$）与随机变量版本（$p(\mathbf{y})>0$）：
\begin{equation}
  \mathbb{P}(A\mid B)=\frac{\mathbb{P}(B\mid A)\,\mathbb{P}(A)}{\mathbb{P}(B)},
  \qquad
  \underbrace{p(\mathbf{x}\mid\mathbf{y})}_{\text{后验}}=\frac{\overbrace{p(\mathbf{y}\mid\mathbf{x})}^{\text{似然}}\ \overbrace{p(\mathbf{x})}^{\text{先验}}}{\underbrace{p(\mathbf{y})}_{\text{证据}}}.
  \label{eq:prob-bayes}
\end{equation}
其中证据（归一化常数）$p(\mathbf{y})=\int p(\mathbf{y}\mid\mathbf{x})\,p(\mathbf{x})\,d\mathbf{x}$ 由全概率公式 \cref{eq:total-prob} 给出，与 $\mathbf{x}$ 无关，故常写成正比形式
\begin{equation}
  p(\mathbf{x}\mid\mathbf{y})\ \propto\ p(\mathbf{y}\mid\mathbf{x})\,p(\mathbf{x}).
  \label{eq:bayes-proportional}
\end{equation}
\end{theorem}
\begin{proof}
由条件概率定义两次写联合：$\mathbb{P}(A\cap B)=\mathbb{P}(A\mid B)\mathbb{P}(B)=\mathbb{P}(B\mid A)\mathbb{P}(A)$，两端除以 $\mathbb{P}(B)$ 即得。随机变量版同理由 $p(\mathbf{x},\mathbf{y})=p(\mathbf{x}\mid\mathbf{y})p(\mathbf{y})=p(\mathbf{y}\mid\mathbf{x})p(\mathbf{x})$ 得证。
\end{proof}

\begin{insight}[贝叶斯定理是“信念的乘法更新”]
读 \cref{eq:bayes-proportional}：\emph{后验} $\propto$ \emph{似然} $\times$ \emph{先验}。先验 $p(\mathbf{x})$ 是看到新数据\emph{前}的信念；似然 $p(\mathbf{y}\mid\mathbf{x})$ 是“若状态为 $\mathbf{x}$，看到此观测的可能性”；二者相乘再归一，即得看到数据\emph{后}的信念 $p(\mathbf{x}\mid\mathbf{y})$。整套贝叶斯滤波（\cref{ch:eskf}）就是把这条更新\emph{递归}地用下去：上一时刻后验当作这一时刻先验，新观测的似然乘上去——“预测—更新”循环的“更新”二字，指的正是这一次贝叶斯乘法。
\end{insight}

\begin{example}[传感器报警的反直觉：为何“报警”远不等于“存在”]\label{ex:bayes-alarm}
承前置自测第 $2$ 题。设障碍物\emph{先验}存在率低，$\mathbb{P}(H)=0.01$；传感器真阳率 $\mathbb{P}(D\mid H)=0.99$、误报率 $\mathbb{P}(D\mid H^{\mathrm c})=0.05$。现报警（$D$），求 $\mathbb{P}(H\mid D)$。由 \cref{eq:prob-bayes} 与全概率：
\[
  \mathbb{P}(H\mid D)=\frac{0.99\times0.01}{0.99\times0.01+0.05\times0.99}
  =\frac{0.0099}{0.0099+0.0495}\approx0.167.
\]
即便传感器“$99\%$ 准”，一次报警下障碍物真存在的概率也只有约 $17\%$！原因是\emph{先验太低}：稀有事件即使被高准确率传感器命中，也会被海量“误报”稀释。这正是\textbf{先验不可忽略}的铁证，也解释了为何机器人不能见一次异常读数就贸然急停——单帧观测须与先验、与多帧融合后才可信。该现象在医学诊断中称“基率谬误”。
\end{example}

\paragraph{相互独立 vs 不相关：一字之差，谬以千里。}
两个常被混为一谈的概念：
\begin{itemize}
  \item \textbf{不相关}：$\operatorname{Cov}[X,Y]=0$，即\emph{无线性}联动；
  \item \textbf{独立}：$p(x,y)=p(x)p(y)$，即\emph{毫无}（任何阶）联动。
\end{itemize}
独立 $\Rightarrow$ 不相关（独立时 $\mathbb{E}[XY]=\mathbb{E}[X]\mathbb{E}[Y]$，故协方差为零），但\textbf{反之不真}。

\begin{example}[不相关却不独立的经典反例]\label{ex:uncorr-not-indep}
设 $X\sim U[-1,1]$（对称），$Y=X^2$。则 $\operatorname{Cov}[X,Y]=\mathbb{E}[X^3]-\mathbb{E}[X]\mathbb{E}[X^2]=0-0=0$（奇函数期望为零），故 $X,Y$ \emph{不相关}。但 $Y$ 完全由 $X$ 决定，二者\emph{极度相关}（非线性），绝不独立——知道 $X=0.5$ 就断定 $Y=0.25$。教训：协方差只看线性关系，\textbf{协方差为零并不意味着“无关”}。在状态估计里把“不相关”当“独立”用，可能漏掉强非线性耦合而过度自信。
\end{example}

\begin{pitfall}[思维陷阱：把“不相关”当“独立”，把对角协方差当“各分量独立”]
工程中常把协方差矩阵设为\emph{对角}以求省事，并默认“各状态分量独立”。\cref{ex:uncorr-not-indep} 表明：对角（即两两不相关）\emph{一般}推不出独立。\textbf{唯一的例外是高斯}——联合高斯下，不相关 $\Longleftrightarrow$ 独立（\cref{cor:gauss-uncorr-indep}），这是高斯的特殊恩典，也是工程上敢用对角高斯的真正依据。一旦分布非高斯，对角协方差只压制了线性耦合，非线性依赖仍在，切勿误以为“解耦”。
\end{pitfall}

\begin{practice}
\begin{enumerate}
  \item 用 \cref{eq:prob-bayes} 重做 \cref{ex:bayes-alarm}，把先验改为 $\mathbb{P}(H)=0.5$，观察后验跃升至何值，体会先验的杠杆作用。
  \item 证明：若 $X,Y$ 独立，则 $\mathbb{E}[XY]=\mathbb{E}[X]\mathbb{E}[Y]$，从而 $\operatorname{Cov}[X,Y]=0$（独立 $\Rightarrow$ 不相关）。
  \item 构造另一个“不相关但不独立”的离散例子（提示：让 $(X,Y)$ 在 $\{(\pm1,0),(0,\pm1)\}$ 上等概率取值）。
  \item （概念）用条件独立 \cref{def:cond-independence} 解释“给定机器人当前位姿，相邻两帧图像的内容条件独立”这一建模假设的含义与边界。
\end{enumerate}
\end{practice}

% ════════════════════════════════════════════════════════════════
\section{高斯分布}
\label{sec:prob-gauss}

\paragraph{动机：为什么是高斯？}
高斯（正态）分布在机器人学中无处不在，原因有三：(i) \emph{解析友好}——线性变换、边缘化、条件化、相乘都\emph{封闭}于高斯，给出闭式滤波；(ii) \emph{两参数足矣}——均值与协方差完全决定分布，滤波只需传播这两个量；(iii) \emph{理论必然}——大量独立小扰动之和趋于高斯（中心极限定理，\cref{sec:prob-limit}），故传感器噪声常近似高斯。本节给出一元与多元高斯，下一节专攻其\emph{条件分布}（滤波核心）。

\begin{definition}[一元高斯]\label{def:gauss-1d}
随机变量 $X$ 服从均值 $\mu$、方差 $\sigma^2$ 的\textbf{高斯分布} $X\sim\mathcal{N}(\mu,\sigma^2)$，若密度为
\begin{equation}
  p(x)=\frac{1}{\sqrt{2\pi}\,\sigma}\exp\!\Big(\!-\frac{(x-\mu)^2}{2\sigma^2}\Big).
  \label{eq:gauss-1d}
\end{equation}
$\mu$ 定中心，$\sigma$ 定宽度；约 $68\%/95\%/99.7\%$ 的概率质量落在 $\mu\pm\sigma/2\sigma/3\sigma$ 内（“$3\sigma$ 法则”，工程上据此设阈值）。
\end{definition}

\begin{definition}[多元高斯]\label{def:gauss-nd}
随机向量 $\mathbf{x}\in\mathbb{R}^n$ 服从\textbf{多元高斯} $\mathbf{x}\sim\mathcal{N}(\boldsymbol{\mu},\boldsymbol{\Sigma})$（$\boldsymbol{\Sigma}$ 对称正定），若
\begin{equation}
  p(\mathbf{x})=\frac{1}{(2\pi)^{n/2}\,(\det\boldsymbol{\Sigma})^{1/2}}\exp\!\Big(\!-\tfrac12(\mathbf{x}-\boldsymbol{\mu})^\top\boldsymbol{\Sigma}^{-1}(\mathbf{x}-\boldsymbol{\mu})\Big).
  \label{eq:gauss-nd}
\end{equation}
指数中的二次型
\begin{equation}
  d_M^2(\mathbf{x})=(\mathbf{x}-\boldsymbol{\mu})^\top\boldsymbol{\Sigma}^{-1}(\mathbf{x}-\boldsymbol{\mu})
  \label{eq:mahalanobis}
\end{equation}
称为\textbf{马氏距离}（Mahalanobis distance）的平方，是“按不确定度加权”的距离。
\end{definition}

\paragraph{马氏距离与协方差椭球的几何。}
$d_M^2(\mathbf{x})=c$（常数）的等高面是一族\emph{椭球}（高斯密度的等高线）。把 $\boldsymbol{\Sigma}$ 特征分解 $\boldsymbol{\Sigma}=\mathbf{U}\boldsymbol{\Lambda}_{\!\mathrm{eig}}\mathbf{U}^\top$（$\mathbf{U}$ 正交、$\boldsymbol{\Lambda}_{\!\mathrm{eig}}=\operatorname{diag}(\lambda_i)$），则椭球\emph{主轴方向}为特征向量 $\mathbf{u}_i$、\emph{半轴长}正比于 $\sqrt{\lambda_i}$。直觉：$\boldsymbol{\Sigma}^{-1}$ 在“方差大”的方向上权重小（容忍大偏差）、在“方差小”的方向上权重大（偏一点就远）。这把“不确定度”可视化成一个椭球——后续协方差传播 $\mathbf{A}\boldsymbol{\Sigma}\mathbf{A}^\top$ 就是这个椭球被线性映射拉伸/旋转。

\begin{figure}[ht]
\centering
\begin{tikzpicture}[scale=1.0,>=Stealth]
  % axes
  \draw[->,gray!60!black] (-2.4,0)--(3.0,0) node[right]{$x_1$};
  \draw[->,gray!60!black] (0,-1.9)--(0,2.4) node[above]{$x_2$};
  % mean
  \coordinate (mu) at (0.6,0.5);
  \fill[main] (mu) circle (1.4pt) node[below right,font=\small] {$\boldsymbol{\mu}$};
  % rotated ellipses: 1,2,3 sigma (rotate 30 deg, axes 0.7 x 1.3)
  \begin{scope}[shift={(mu)}, rotate=30]
    \draw[second!75!black, thick] (0,0) ellipse (0.7 and 1.3);
    \draw[second!55!black] (0,0) ellipse (1.4 and 2.6);
    % principal axes
    \draw[->,third!70!black,thick] (0,0)--(0.7,0) node[midway,below,font=\scriptsize]{$\sqrt{\lambda_1}\,\mathbf{u}_1$};
    \draw[->,third!70!black,thick] (0,0)--(0,1.3) node[midway,left,font=\scriptsize]{$\sqrt{\lambda_2}\,\mathbf{u}_2$};
  \end{scope}
  \node[second!60!black,font=\small] at (2.55,1.95) {$d_M^2=c$};
\end{tikzpicture}
\caption{二维高斯的协方差椭球：等马氏距离面 $d_M^2(\mathbf{x})=c$ 是以 $\boldsymbol{\mu}$ 为心的椭球，主轴沿 $\boldsymbol{\Sigma}$ 的特征向量 $\mathbf{u}_i$、半轴正比于 $\sqrt{\lambda_i}$。内外两椭圈示意 $1\sigma$ 与 $2\sigma$。}
\label{fig:cov-ellipse}
\end{figure}

\begin{theorem}[高斯在线性变换与边缘化下封闭]\label{thm:gauss-closed}
设 $\mathbf{x}\sim\mathcal{N}(\boldsymbol{\mu},\boldsymbol{\Sigma})$。
\begin{enumerate}
  \item \textbf{线性变换}：$\mathbf{y}=\mathbf{A}\mathbf{x}+\mathbf{b}$（$\mathbf{A}$ 行满秩）服从 $\mathbf{y}\sim\mathcal{N}(\mathbf{A}\boldsymbol{\mu}+\mathbf{b},\ \mathbf{A}\boldsymbol{\Sigma}\mathbf{A}^\top)$；
  \item \textbf{边缘化}：把 $\mathbf{x}=[\mathbf{x}_a^\top,\mathbf{x}_b^\top]^\top$ 分块，则边缘 $\mathbf{x}_a\sim\mathcal{N}(\boldsymbol{\mu}_a,\boldsymbol{\Sigma}_{aa})$——\emph{直接取对应子块即可}。
\end{enumerate}
\end{theorem}
\begin{proof}[证明骨架]
线性变换：均值与协方差由 \cref{eq:affine-prop} 给出；要证结果\emph{仍是高斯}，可用矩母函数 $M_{\mathbf{x}}(\mathbf{t})=\mathbb{E}[e^{\mathbf{t}^\top\mathbf{x}}]=\exp(\mathbf{t}^\top\boldsymbol{\mu}+\tfrac12\mathbf{t}^\top\boldsymbol{\Sigma}\mathbf{t})$，则 $M_{\mathbf{y}}(\mathbf{t})=e^{\mathbf{t}^\top\mathbf{b}}M_{\mathbf{x}}(\mathbf{A}^\top\mathbf{t})=\exp(\mathbf{t}^\top(\mathbf{A}\boldsymbol{\mu}+\mathbf{b})+\tfrac12\mathbf{t}^\top\mathbf{A}\boldsymbol{\Sigma}\mathbf{A}^\top\mathbf{t})$，正是高斯的矩母函数，由其唯一性知 $\mathbf{y}$ 高斯。边缘化是线性变换 $\mathbf{x}_a=[\mathbf{I}\ \mathbf{0}]\mathbf{x}$ 的特例，代入即得 $\boldsymbol{\Sigma}_{aa}$。
\end{proof}

\begin{corollary}[联合高斯下：不相关 $\Longleftrightarrow$ 独立]\label{cor:gauss-uncorr-indep}
若 $(\mathbf{x}_a,\mathbf{x}_b)$ \emph{联合}高斯，则 $\mathbf{x}_a,\mathbf{x}_b$ 独立当且仅当 $\boldsymbol{\Sigma}_{ab}=\operatorname{Cov}[\mathbf{x}_a,\mathbf{x}_b]=\mathbf{0}$。
\end{corollary}
\begin{proof}
$\boldsymbol{\Sigma}_{ab}=\mathbf{0}$ 时 $\boldsymbol{\Sigma}$ 块对角，$\boldsymbol{\Sigma}^{-1}$ 亦块对角，\cref{eq:gauss-nd} 的指数二次型按块分裂，密度 $p(\mathbf{x}_a,\mathbf{x}_b)=p(\mathbf{x}_a)p(\mathbf{x}_b)$，即独立。反向由 \cref{ex:uncorr-not-indep} 之逆——独立总蕴含不相关。\textbf{强调}：此等价\emph{仅}对联合高斯成立（\cref{ex:uncorr-not-indep} 的非高斯反例不矛盾，因其非联合高斯）。
\end{proof}

\begin{pitfall}[概念误区：以为“边缘高斯”就等于“联合高斯”]
\cref{cor:gauss-uncorr-indep} 的前提是\emph{联合}高斯。两个分量\emph{各自}（边缘）高斯，\emph{并不}保证它们\emph{联合}高斯——存在边缘都高斯、联合却非高斯的反例。因此“$X$ 高斯、$Y$ 高斯”不足以推出“$(X,Y)$ 高斯”，更不能据边缘的不相关推独立。状态估计里默认的“联合高斯”是个\emph{建模假设}，需由线性高斯模型（线性动力学 $+$ 高斯噪声）来\emph{保证}（\cref{ch:eskf}），而非天然成立。
\end{pitfall}

\begin{practice}
\begin{enumerate}
  \item 由 \cref{eq:gauss-nd} 验证归一化：$\int_{\mathbb{R}^n}p(\mathbf{x})\,d\mathbf{x}=1$（提示：作白化变换 $\mathbf{z}=\boldsymbol{\Sigma}^{-1/2}(\mathbf{x}-\boldsymbol{\mu})$ 化为标准高斯，$d\mathbf{x}=(\det\boldsymbol{\Sigma})^{1/2}d\mathbf{z}$）。
  \item 证明马氏距离平方 $d_M^2(\mathbf{x})\sim\chi^2_n$（卡方分布，自由度 $n$），从而 $\mathbb{E}[d_M^2]=n$——这是滤波一致性检验（NEES/NIS）的理论依据（\cref{ch:eskf}）。
  \item 取 $\boldsymbol{\Sigma}=\big[\begin{smallmatrix}4&0\\0&1\end{smallmatrix}\big]$ 与 $\big[\begin{smallmatrix}2&1\\1&2\end{smallmatrix}\big]$，分别画出 $d_M^2=1$ 的椭圆，标出主轴方向与半轴长。
\end{enumerate}
\end{practice}

% ════════════════════════════════════════════════════════════════
\section{多元高斯的条件分布：滤波核心}
\label{sec:prob-gauss-cond}

\paragraph{动机：这是卡尔曼滤波的数学心脏。}
状态估计的基本动作是：把状态 $\mathbf{x}$ 与观测 $\mathbf{y}$ 视为\emph{联合}高斯，观测到具体的 $\mathbf{y}$ 后，求 $\mathbf{x}$ 的\emph{条件}分布 $p(\mathbf{x}\mid\mathbf{y})$。若这仍是高斯、且其均值/协方差有闭式，则“由观测更新状态”就是一次矩阵运算——这正是卡尔曼增益的来源。本节给出并证明这条公式，它是 \cref{ch:eskf} 卡尔曼更新步、\cref{ch:slam_est} 高斯条件化的共同源头。

\begin{theorem}[联合高斯的条件分布]\label{thm:gauss-cond}
设 $\begin{bmatrix}\mathbf{x}\\\mathbf{y}\end{bmatrix}\sim\mathcal{N}\!\left(\begin{bmatrix}\boldsymbol{\mu}_x\\\boldsymbol{\mu}_y\end{bmatrix},\ \begin{bmatrix}\boldsymbol{\Sigma}_{xx}&\boldsymbol{\Sigma}_{xy}\\\boldsymbol{\Sigma}_{yx}&\boldsymbol{\Sigma}_{yy}\end{bmatrix}\right)$，则给定 $\mathbf{y}$ 的条件分布仍为高斯 $\mathbf{x}\mid\mathbf{y}\sim\mathcal{N}(\boldsymbol{\mu}_{x\mid y},\boldsymbol{\Sigma}_{x\mid y})$，且
\begin{align}
  \boldsymbol{\mu}_{x\mid y}&=\boldsymbol{\mu}_x+\boldsymbol{\Sigma}_{xy}\boldsymbol{\Sigma}_{yy}^{-1}(\mathbf{y}-\boldsymbol{\mu}_y),\label{eq:cond-mean}\\
  \boldsymbol{\Sigma}_{x\mid y}&=\boldsymbol{\Sigma}_{xx}-\boldsymbol{\Sigma}_{xy}\boldsymbol{\Sigma}_{yy}^{-1}\boldsymbol{\Sigma}_{yx}.\label{eq:cond-cov}
\end{align}
其中 \cref{eq:cond-cov} 的右端正是分块矩阵 $\boldsymbol{\Sigma}$ 关于 $\boldsymbol{\Sigma}_{yy}$ 的\textbf{舒尔补}（Schur complement）。
\end{theorem}
\begin{proof}[证明骨架：分块求逆 $+$ 配方]
记联合精度矩阵 $\boldsymbol{\Lambda}=\boldsymbol{\Sigma}^{-1}=\big[\begin{smallmatrix}\boldsymbol{\Lambda}_{xx}&\boldsymbol{\Lambda}_{xy}\\\boldsymbol{\Lambda}_{yx}&\boldsymbol{\Lambda}_{yy}\end{smallmatrix}\big]$。联合密度的指数（去掉常数）为 $-\tfrac12\big[(\mathbf{x}-\boldsymbol{\mu}_x)^\top,(\mathbf{y}-\boldsymbol{\mu}_y)^\top\big]\boldsymbol{\Lambda}\big[\cdots\big]$。把 $\mathbf{y}$ 视为常数、只收集关于 $\mathbf{x}$ 的二次与一次项，配方成 $-\tfrac12(\mathbf{x}-\mathbf{m})^\top\boldsymbol{\Lambda}_{xx}(\mathbf{x}-\mathbf{m})+\text{const}$ 的形式，直接读出
\[
  \boldsymbol{\Sigma}_{x\mid y}=\boldsymbol{\Lambda}_{xx}^{-1},\qquad
  \mathbf{m}=\boldsymbol{\mu}_{x\mid y}=\boldsymbol{\mu}_x-\boldsymbol{\Lambda}_{xx}^{-1}\boldsymbol{\Lambda}_{xy}(\mathbf{y}-\boldsymbol{\mu}_y).
\]
再用分块求逆公式把精度块 $\boldsymbol{\Lambda}_{xx},\boldsymbol{\Lambda}_{xy}$ 用协方差块表出：$\boldsymbol{\Lambda}_{xx}=(\boldsymbol{\Sigma}_{xx}-\boldsymbol{\Sigma}_{xy}\boldsymbol{\Sigma}_{yy}^{-1}\boldsymbol{\Sigma}_{yx})^{-1}$（即舒尔补之逆），$\boldsymbol{\Lambda}_{xx}^{-1}\boldsymbol{\Lambda}_{xy}=-\boldsymbol{\Sigma}_{xy}\boldsymbol{\Sigma}_{yy}^{-1}$，代入即得 \cref{eq:cond-mean,eq:cond-cov}。配方与分块求逆的完整代数在 \cref{ch:slam_est} 的信息形式一节展开，此处给出骨架以见其必然。
\end{proof}

\begin{insight}[条件均值是“先验 $+$ 增益 $\times$ 新息”——卡尔曼增益现身]
凝视 \cref{eq:cond-mean}：条件均值 $=$ 先验均值 $\boldsymbol{\mu}_x$ \emph{加上}一个修正项 $\boldsymbol{\Sigma}_{xy}\boldsymbol{\Sigma}_{yy}^{-1}(\mathbf{y}-\boldsymbol{\mu}_y)$。其中 $(\mathbf{y}-\boldsymbol{\mu}_y)$ 是\textbf{新息}（观测减去其先验预测），而矩阵 $\mathbf{K}=\boldsymbol{\Sigma}_{xy}\boldsymbol{\Sigma}_{yy}^{-1}$ 正是\textbf{卡尔曼增益}的雏形——它衡量“状态与观测的相关 $\boldsymbol{\Sigma}_{xy}$”相对于“观测自身的不确定 $\boldsymbol{\Sigma}_{yy}$”，决定新息该被多大程度地采信。\cref{eq:cond-cov} 则说：观测\emph{总是}减小不确定度（$\boldsymbol{\Sigma}_{x\mid y}\preceq\boldsymbol{\Sigma}_{xx}$，因减去半正定的舒尔补项）。把观测模型 $\mathbf{y}=\mathbf{H}\mathbf{x}+\text{噪声}$ 代入（于是 $\boldsymbol{\Sigma}_{xy}=\boldsymbol{\Sigma}_{xx}\mathbf{H}^\top$、$\boldsymbol{\Sigma}_{yy}=\mathbf{H}\boldsymbol{\Sigma}_{xx}\mathbf{H}^\top+\mathbf{R}$），\cref{eq:cond-mean,eq:cond-cov} 立刻\emph{原样变成}卡尔曼滤波的更新公式（\cref{ch:eskf}）。本定理就是卡尔曼滤波“为什么长那样”的全部答案。
\end{insight}

\begin{theorem}[两个高斯之积仍是（未归一）高斯]\label{thm:gauss-product}
两个关于同一变量 $\mathbf{x}$ 的高斯密度之积仍正比于一个高斯：$\mathcal{N}(\mathbf{x};\boldsymbol{\mu}_1,\boldsymbol{\Sigma}_1)\cdot\mathcal{N}(\mathbf{x};\boldsymbol{\mu}_2,\boldsymbol{\Sigma}_2)\propto\mathcal{N}(\mathbf{x};\boldsymbol{\mu},\boldsymbol{\Sigma})$，其中（以\emph{信息形式}最简洁）
\begin{equation}
  \boldsymbol{\Sigma}^{-1}=\boldsymbol{\Sigma}_1^{-1}+\boldsymbol{\Sigma}_2^{-1},
  \qquad
  \boldsymbol{\Sigma}^{-1}\boldsymbol{\mu}=\boldsymbol{\Sigma}_1^{-1}\boldsymbol{\mu}_1+\boldsymbol{\Sigma}_2^{-1}\boldsymbol{\mu}_2.
  \label{eq:gauss-product}
\end{equation}
\end{theorem}
\begin{proof}
两高斯指数相加，$-\tfrac12(\mathbf{x}-\boldsymbol{\mu}_1)^\top\boldsymbol{\Sigma}_1^{-1}(\mathbf{x}-\boldsymbol{\mu}_1)-\tfrac12(\mathbf{x}-\boldsymbol{\mu}_2)^\top\boldsymbol{\Sigma}_2^{-1}(\mathbf{x}-\boldsymbol{\mu}_2)$，按 $\mathbf{x}$ 配方：二次项系数给 $\boldsymbol{\Sigma}^{-1}=\boldsymbol{\Sigma}_1^{-1}+\boldsymbol{\Sigma}_2^{-1}$，一次项给 $\boldsymbol{\Sigma}^{-1}\boldsymbol{\mu}=\boldsymbol{\Sigma}_1^{-1}\boldsymbol{\mu}_1+\boldsymbol{\Sigma}_2^{-1}\boldsymbol{\mu}_2$，余项并入归一化常数。
\end{proof}

\begin{insight}[“信息相加”——为何后端偏爱信息矩阵]
\cref{eq:gauss-product} 揭示一条优雅事实：高斯\emph{相乘}（即贝叶斯融合两条独立证据）在\textbf{信息（精度）形式}下就是\emph{相加} $\boldsymbol{\Sigma}^{-1}=\boldsymbol{\Sigma}_1^{-1}+\boldsymbol{\Sigma}_2^{-1}$。这把“融合 $N$ 个独立观测”从一串矩阵乘法降为一次求和，且天然可并行、可任意顺序。SLAM 后端、因子图、信息滤波之所以钟爱信息矩阵 $\boldsymbol{\Lambda}=\boldsymbol{\Sigma}^{-1}$（\cref{ch:slam_est}），根子就在这里——观测的“信息”是可加的货币。结合 \cref{thm:bayes}：贝叶斯更新 $=$ 先验信息 $+$ 似然信息。
\end{insight}

\begin{table}[H]\centering\small
\caption{协方差形式与信息形式的互补（高斯下）}
\begin{tabular}{lll}
\toprule
操作 & 协方差形式 $(\boldsymbol{\mu},\boldsymbol{\Sigma})$ & 信息形式 $(\boldsymbol{\eta}=\boldsymbol{\Lambda}\boldsymbol{\mu},\ \boldsymbol{\Lambda}=\boldsymbol{\Sigma}^{-1})$ \\
\midrule
边缘化（积掉变量） & 易：取子块 $\boldsymbol{\Sigma}_{aa}$ & 难：需舒尔补 \\
条件化（固定观测） & 难：需舒尔补（\cref{eq:cond-cov}） & 易：取子块 $\boldsymbol{\Lambda}_{aa}$ \\
融合独立观测（相乘） & 难：串行增益 & 易：相加（\cref{eq:gauss-product}） \\
线性变换 $\mathbf{A}\mathbf{x}$ & 易：$\mathbf{A}\boldsymbol{\Sigma}\mathbf{A}^\top$ & 难 \\
“一无所知”表示 & 不可（$\boldsymbol{\Sigma}\to\infty$） & 易（$\boldsymbol{\Lambda}\to\mathbf{0}$） \\
\bottomrule
\end{tabular}
\end{table}

\begin{practice}
\begin{enumerate}
  \item 一维特例：$X\sim\mathcal{N}(\mu_x,\sigma_x^2)$，观测 $Y=X+v$、$v\sim\mathcal{N}(0,\sigma_v^2)$ 独立。由 \cref{thm:gauss-cond} 算 $X\mid Y=y$ 的均值与方差，验证均值 $=\mu_x+\frac{\sigma_x^2}{\sigma_x^2+\sigma_v^2}(y-\mu_x)$，并指出增益 $\frac{\sigma_x^2}{\sigma_x^2+\sigma_v^2}\in(0,1)$ 的含义。
  \item 由 \cref{eq:gauss-product} 推一维高斯融合：$\sigma^{-2}=\sigma_1^{-2}+\sigma_2^{-2}$，$\mu=\sigma^2(\sigma_1^{-2}\mu_1+\sigma_2^{-2}\mu_2)$，说明“融合后必更确定”（$\sigma<\min(\sigma_1,\sigma_2)$）。
  \item 把观测模型 $\mathbf{y}=\mathbf{H}\mathbf{x}+\mathbf{v}$（$\mathbf{v}\sim\mathcal{N}(\mathbf{0},\mathbf{R})$，$\mathbf{x}\sim\mathcal{N}(\boldsymbol{\mu}_x,\boldsymbol{\Sigma}_{xx})$）代入 \cref{eq:cond-mean,eq:cond-cov}，导出卡尔曼增益 $\mathbf{K}=\boldsymbol{\Sigma}_{xx}\mathbf{H}^\top(\mathbf{H}\boldsymbol{\Sigma}_{xx}\mathbf{H}^\top+\mathbf{R})^{-1}$ 与更新 $\boldsymbol{\Sigma}_{x\mid y}=(\mathbf{I}-\mathbf{K}\mathbf{H})\boldsymbol{\Sigma}_{xx}$。
\end{enumerate}
\end{practice}

% ════════════════════════════════════════════════════════════════
\section{大数定律与中心极限定理}
\label{sec:prob-limit}

\paragraph{动机：两条让概率“可工程化”的极限律。}
两条渐近定律支撑着机器人学的两大实践：\emph{大数定律}（LLN）说“样本平均收敛到真期望”，是\textbf{蒙特卡洛}（粒子滤波、采样估计）的理论保证；\emph{中心极限定理}（CLT）说“大量独立小量之和近似高斯”，是\textbf{“到处假设高斯”}的底气。

\begin{theorem}[（弱）大数定律]\label{thm:lln}
设 $X_1,X_2,\dots$ 独立同分布（i.i.d.），期望 $\mathbb{E}[X_i]=\mu$ 有限。则样本均值 $\bar X_n=\frac1n\sum_{i=1}^n X_i$ 依概率收敛到 $\mu$：对任意 $\varepsilon>0$，$\lim_{n\to\infty}\mathbb{P}(|\bar X_n-\mu|>\varepsilon)=0$。
\end{theorem}
\begin{proof}[证明骨架（设方差 $\sigma^2$ 有限）]
由独立性 $\operatorname{Var}[\bar X_n]=\sigma^2/n$（\cref{eq:cov-additive} 加 \cref{eq:variance} 的缩放律）。由切比雪夫不等式 $\mathbb{P}(|\bar X_n-\mu|>\varepsilon)\le\operatorname{Var}[\bar X_n]/\varepsilon^2=\sigma^2/(n\varepsilon^2)\to0$。（去掉有限方差假设的版本需更细的截断技巧。）
\end{proof}

\begin{theorem}[中心极限定理]\label{thm:clt}
设 $X_1,\dots,X_n$ i.i.d.，期望 $\mu$、方差 $\sigma^2\in(0,\infty)$。则标准化样本均值依分布收敛到标准高斯：
\begin{equation}
  \frac{\bar X_n-\mu}{\sigma/\sqrt{n}}\ \xrightarrow{\ d\ }\ \mathcal{N}(0,1),
  \quad\text{即}\quad
  \sum_{i=1}^n X_i\ \approx\ \mathcal{N}(n\mu,\,n\sigma^2)\ (n\text{ 大}).
  \label{eq:clt}
\end{equation}
\end{theorem}
\begin{proof}[证明骨架（特征函数法）]
设 $Y_i=(X_i-\mu)/\sigma$（零均值单位方差），其特征函数 $\varphi_Y(t)=1-\tfrac{t^2}{2}+o(t^2)$。$\frac1{\sqrt n}\sum Y_i$ 的特征函数为 $\varphi_Y(t/\sqrt n)^n=\big(1-\tfrac{t^2}{2n}+o(1/n)\big)^n\to e^{-t^2/2}$，正是 $\mathcal{N}(0,1)$ 的特征函数；由连续性定理得依分布收敛。
\end{proof}

\begin{insight}[CLT 是“高斯无处不在”的原因，但要看清它的边界]
为何传感器噪声常被建模为高斯？因为单次读数的误差往往是\emph{许多}独立微扰（热噪声、量化、机械振动、电路干扰……）之\emph{和}，由 CLT 其和近似高斯——这\emph{解释}了高斯假设的合理性，而非凭空假设。但 CLT 有前提与边界：(i) 各项须\emph{独立}且方差有限——若存在重尾（如某些异常/野值机制），和未必趋高斯，这正是\emph{鲁棒估计}要处理的情形；(ii) 收敛是对\emph{中心}的，\emph{尾部}逼近慢——故 $3\sigma$ 之外的稀有大偏差用高斯估计常不准，外点检测（如卡方门限，\cref{ch:eskf}）由此而生。把“可以用高斯”与“误差一定服从高斯”分清，是工程清醒的标志。
\end{insight}

\begin{example}[蒙特卡洛估计的收敛速率]\label{ex:monte-carlo}
要估计期望 $\mathbb{E}[g(X)]$ 而积分难算时，抽 $n$ 个独立样本 $x_i\sim p$，用样本平均 $\hat I_n=\frac1n\sum g(x_i)$。由 LLN，$\hat I_n\to\mathbb{E}[g(X)]$；由 CLT，误差 $\hat I_n-\mathbb{E}[g(X)]\approx\mathcal{N}(0,\operatorname{Var}[g(X)]/n)$，即标准误差按 $1/\sqrt n$ 下降——\emph{与维度无关}。这正是粒子滤波、采样式规划（如 \cref{ch:slam_est} 之外的采样估计）敢用蒙特卡洛对抗“维度灾难”的根据：精度翻倍要 $4$ 倍样本，虽慢但不随维数爆炸。
\end{example}

\begin{practice}
\begin{enumerate}
  \item 用切比雪夫不等式补全 \cref{thm:lln} 的证明，并据此说明：要把 $|\bar X_n-\mu|>\varepsilon$ 的概率压到 $\delta$ 以下，样本数 $n$ 至少要多大（用 $\sigma,\varepsilon,\delta$ 表示）。
  \item （数值实验思路）取一个明显非高斯的分布（如均匀分布），数值地抽 $n=1,2,5,30$ 个样本求均值、重复多次画直方图，观察 $n$ 增大时直方图如何向钟形靠拢——亲手验证 CLT。
  \item 由 \cref{ex:monte-carlo} 说明：蒙特卡洛误差 $\propto1/\sqrt n$ 而与维度无关，对比确定性数值积分（网格法）误差随维度指数恶化，解释高维下采样法的优势。
\end{enumerate}
\end{practice}

% ════════════════════════════════════════════════════════════════
\section{KL 散度、熵与信息量}
\label{sec:prob-kl}

\paragraph{动机：给“分布之间的差异”和“不确定度”一把尺子。}
机器人学常需\emph{比较两个分布}（真实后验 vs 近似后验，用于变分推断）、或\emph{量化一个分布的不确定度}（熵，用于主动感知“下一步去哪最能消除歧义”）、或\emph{度量两量的关联}（互信息，用于特征选择与信息增益）。KL 散度、熵、互信息正是这三把尺子，本节给出定义与机器人学最常用的高斯情形。

\begin{definition}[（微分）熵]\label{def:entropy}
离散分布 $p$ 的\textbf{香农熵}与连续分布的\textbf{微分熵}：
\begin{equation}
  H(p)=-\sum_k p(x_k)\log p(x_k),
  \qquad
  h(p)=-\int p(\mathbf{x})\log p(\mathbf{x})\,d\mathbf{x}=-\mathbb{E}_p[\log p(\mathbf{x})].
  \label{eq:entropy}
\end{equation}
熵度量分布的“不确定度/平均信息量”：越平坦（越不确定）熵越大，越集中越小。
\end{definition}

\begin{definition}[KL 散度（相对熵）]\label{def:kl}
分布 $p$ 相对 $q$ 的 \textbf{KL 散度}：
\begin{equation}
  \mathrm{KL}(p\,\|\,q)=\int p(\mathbf{x})\log\frac{p(\mathbf{x})}{q(\mathbf{x})}\,d\mathbf{x}=\mathbb{E}_p\!\Big[\log\frac{p}{q}\Big].
  \label{eq:kl}
\end{equation}
它度量“用 $q$ 近似 $p$ 的低效程度”。
\end{definition}

\begin{theorem}[Gibbs 不等式：KL 非负]\label{thm:kl-nonneg}
$\mathrm{KL}(p\,\|\,q)\ge0$，等号当且仅当 $p=q$（几乎处处）。
\end{theorem}
\begin{proof}
由 $\log$ 凹性与 Jensen 不等式：$-\mathrm{KL}(p\|q)=\mathbb{E}_p[\log\frac{q}{p}]\le\log\mathbb{E}_p[\frac{q}{p}]=\log\int q\,d\mathbf{x}=\log1=0$，故 $\mathrm{KL}\ge0$；$\log$ 严格凹，等号仅当 $q/p$ 几乎处处为常数，结合两者均归一即 $p=q$。
\end{proof}

\begin{pitfall}[概念误区：把 KL 散度当“距离”]
KL 记号像距离，但它\textbf{不是度量}：一般\emph{不对称} $\mathrm{KL}(p\|q)\ne\mathrm{KL}(q\|p)$，也\emph{不满足}三角不等式。方向有实义：变分推断里最小化 $\mathrm{KL}(q\|p)$（用近似 $q$ 去“贴”真后验 $p$）会得到\emph{众数寻求}（mode-seeking）、偏窄的近似；最小化 $\mathrm{KL}(p\|q)$ 则\emph{覆盖}所有模式、偏宽。选错方向，近似分布的“胖瘦”会系统性偏差。需要真正的对称距离时，用 Jensen--Shannon 散度或 Wasserstein 距离。
\end{pitfall}

\begin{theorem}[高斯的熵与 KL（机器人学最常用）]\label{thm:gauss-info}
对 $\mathbf{x}\sim\mathcal{N}(\boldsymbol{\mu},\boldsymbol{\Sigma})$（$n$ 维），微分熵
\begin{equation}
  h(\mathcal{N})=\tfrac12\log\big((2\pi e)^n\det\boldsymbol{\Sigma}\big)=\tfrac{n}{2}\log(2\pi e)+\tfrac12\log\det\boldsymbol{\Sigma}.
  \label{eq:prob-gauss-entropy}
\end{equation}
两高斯 $p_i=\mathcal{N}(\boldsymbol{\mu}_i,\boldsymbol{\Sigma}_i)$ 间的 KL 散度
\begin{equation}
  \mathrm{KL}(p_2\,\|\,p_1)=\tfrac12\Big[(\boldsymbol{\mu}_2-\boldsymbol{\mu}_1)^\top\boldsymbol{\Sigma}_1^{-1}(\boldsymbol{\mu}_2-\boldsymbol{\mu}_1)+\operatorname{tr}(\boldsymbol{\Sigma}_1^{-1}\boldsymbol{\Sigma}_2)-n+\log\frac{\det\boldsymbol{\Sigma}_1}{\det\boldsymbol{\Sigma}_2}\Big].
  \label{eq:prob-gauss-kl}
\end{equation}
\end{theorem}
\begin{proof}[证明骨架]
熵：把 $-\log p(\mathbf{x})=\tfrac12\log((2\pi)^n\det\boldsymbol{\Sigma})+\tfrac12(\mathbf{x}-\boldsymbol{\mu})^\top\boldsymbol{\Sigma}^{-1}(\mathbf{x}-\boldsymbol{\mu})$ 取期望，用 $\mathbb{E}[(\mathbf{x}-\boldsymbol{\mu})^\top\boldsymbol{\Sigma}^{-1}(\mathbf{x}-\boldsymbol{\mu})]=\operatorname{tr}(\boldsymbol{\Sigma}^{-1}\mathbb{E}[(\mathbf{x}-\boldsymbol{\mu})(\mathbf{x}-\boldsymbol{\mu})^\top])=\operatorname{tr}(\mathbf{I})=n$（迹技巧）即得 \cref{eq:prob-gauss-entropy}。KL：把 \cref{eq:kl} 写成 $\mathbb{E}_{p_2}[\log p_2-\log p_1]$，逐项代入两高斯的对数密度、再用同一迹技巧算二次型期望，整理即得 \cref{eq:prob-gauss-kl}。
\end{proof}

\begin{insight}[熵 $\propto\log\det\boldsymbol{\Sigma}$：把“减少不确定度”变成可优化目标]
\cref{eq:prob-gauss-entropy} 表明高斯的不确定度由 $\log\det\boldsymbol{\Sigma}$ 主导，而 $\sqrt{\det\boldsymbol{\Sigma}}$ 正比于协方差椭球（\cref{fig:cov-ellipse}）的体积。于是“信息增益”$=$ 观测前后熵之差 $=\tfrac12\log\frac{\det\boldsymbol{\Sigma}_{\text{先}}}{\det\boldsymbol{\Sigma}_{\text{后}}}$，把“这次观测/这条路径能消除多少不确定”变成一个\emph{可计算、可最大化}的标量——这是主动感知、下一最佳视点（NBV）、信息式探索的目标函数雏形。\cref{eq:prob-gauss-kl} 的二次型项 $(\boldsymbol{\mu}_2-\boldsymbol{\mu}_1)^\top\boldsymbol{\Sigma}_1^{-1}(\cdot)$ 又恰是马氏距离（\cref{eq:mahalanobis}）——“取负对数把高斯变成马氏二次型”这条桥（\cref{ch:slam_est} 的最小二乘由来）在此再次显形。
\end{insight}

\begin{definition}[互信息]\label{def:mutual-info}
随机向量 $\mathbf{x},\mathbf{y}$ 的\textbf{互信息}为联合分布与边缘之积的 KL：
\begin{equation}
  I(\mathbf{x};\mathbf{y})=\mathrm{KL}\big(p(\mathbf{x},\mathbf{y})\,\big\|\,p(\mathbf{x})p(\mathbf{y})\big)=H(\mathbf{x})-H(\mathbf{x}\mid\mathbf{y})\ \ge0.
  \label{eq:prob-mutual-info}
\end{equation}
它度量“知道 $\mathbf{y}$ 平均能消除多少关于 $\mathbf{x}$ 的不确定度”；$I=0$ 当且仅当 $\mathbf{x},\mathbf{y}$ 独立。
\end{definition}
\noindent 互信息 \cref{eq:prob-mutual-info} 是“信息增益”的严格名字：主动 SLAM 中“去哪测最划算”就是\emph{最大化}候选动作与地图的互信息。它由 KL（\cref{def:kl}）与熵（\cref{def:entropy}）自然导出，三者构成信息论的小三角。

\begin{practice}
\begin{enumerate}
  \item 由 \cref{eq:prob-gauss-kl} 取 $\boldsymbol{\Sigma}_1=\boldsymbol{\Sigma}_2=\boldsymbol{\Sigma}$，验证两等协方差高斯的 KL 退化为半个马氏距离 $\tfrac12(\boldsymbol{\mu}_2-\boldsymbol{\mu}_1)^\top\boldsymbol{\Sigma}^{-1}(\boldsymbol{\mu}_2-\boldsymbol{\mu}_1)$。
  \item 一维：由 \cref{eq:prob-gauss-entropy} 得 $h(\mathcal{N}(\mu,\sigma^2))=\tfrac12\log(2\pi e\sigma^2)$，说明熵只依赖 $\sigma$ 不依赖 $\mu$（不确定度与中心位置无关），并据此论证“信息增益与均值无关、只看协方差收缩”。
  \item 用 \cref{eq:prob-mutual-info} 与 \cref{cor:gauss-uncorr-indep} 证明：联合高斯下 $I(\mathbf{x};\mathbf{y})=0\Leftrightarrow\boldsymbol{\Sigma}_{xy}=\mathbf{0}\Leftrightarrow$ 独立。
\end{enumerate}
\end{practice}

% ════════════════════════════════════════════════════════════════
\section*{本章常见误解汇总}
\begin{table}[H]\centering\small
\begin{tabular}{p{0.46\textwidth} p{0.46\textwidth}}
\toprule
\textbf{常见误解} & \textbf{正确理解} \\
\midrule
概率就是长期频率，单次事件无概率 & 公理不限解释；机器人学用“信念程度”同样合法（\cref{def:prob-space}）\\
连续分布的密度值 $p(x)$ 就是概率 & 是密度，可 $>1$；区间积分才是概率（\cref{def:cdf-pdf}）\\
不相关 $=$ 独立 & 独立 $\Rightarrow$ 不相关，反之不真；仅高斯下等价（\cref{cor:gauss-uncorr-indep}）\\
对角协方差 $=$ 各分量独立 & 仅“线性解耦”；非高斯仍可有非线性依赖（\cref{ex:uncorr-not-indep}）\\
报警 $99\%$ 准 $\Rightarrow$ 物体九成存在 & 须乘先验；低先验下后验可低至 $17\%$（\cref{ex:bayes-alarm}）\\
边缘都高斯 $\Rightarrow$ 联合高斯 & 不然；联合高斯是更强的建模假设（\cref{sec:prob-gauss} 陷阱）\\
观测可能增大不确定度 & 高斯条件化\emph{必}减不确定度（\cref{eq:cond-cov} 减半正定项）\\
卡尔曼增益是经验调参 & 是 $\boldsymbol{\Sigma}_{xy}\boldsymbol{\Sigma}_{yy}^{-1}$ 的精确博弈（\cref{thm:gauss-cond}）\\
误差“一定”服从高斯 & CLT 只保证中心近似、独立有限方差；尾部/重尾要鲁棒处理（\cref{thm:clt}）\\
KL 散度是分布间的距离 & 不对称、不满足三角；选错方向致近似胖瘦偏差（\cref{def:kl} 陷阱）\\
熵依赖均值位置 & 高斯熵只依赖 $\det\boldsymbol{\Sigma}$，与 $\boldsymbol{\mu}$ 无关（\cref{eq:prob-gauss-entropy}）\\
\bottomrule
\end{tabular}
\end{table}

% ════════════════════════════════════════════════════════════════
\section*{本章小结}

\begin{table}[H]\centering\small
\caption{符号表}
\begin{tabular}{lll}
\toprule
符号 & 含义 & 首次出现 \\
\midrule
$(\Omega,\mathcal{F},\mathbb{P})$ & 概率空间：样本空间 / 事件域 / 测度 & \cref{def:prob-space} \\
$X,Y$ / $\mathbf{x},\mathbf{y}$ & 随机变量 / 随机向量 & \cref{def:random-variable} \\
$F_X,\ p(x)$ & 分布函数 / 密度（或质量）函数 & \cref{def:cdf-pdf} \\
$p(\mathbf{x},\mathbf{y}),p(\mathbf{x}),p(\mathbf{x}\mid\mathbf{y})$ & 联合 / 边缘 / 条件分布 & \cref{eq:marginal,eq:conditional} \\
$\mathbb{E}[\cdot],\boldsymbol{\mu}$ & 期望 / 均值向量 & \cref{def:expectation} \\
$\operatorname{Var},\operatorname{Cov},\boldsymbol{\Sigma}$ & 方差 / 协方差 / 协方差矩阵 & \cref{def:covariance} \\
$\boldsymbol{\Lambda}=\boldsymbol{\Sigma}^{-1}$ & 信息（精度）矩阵 & \cref{thm:gauss-product} \\
$\mathcal{N}(\boldsymbol{\mu},\boldsymbol{\Sigma})$ & 多元高斯分布 & \cref{def:gauss-nd} \\
$d_M^2=(\mathbf{x}-\boldsymbol{\mu})^\top\boldsymbol{\Sigma}^{-1}(\mathbf{x}-\boldsymbol{\mu})$ & 马氏距离平方 & \cref{eq:mahalanobis} \\
$\boldsymbol{\mu}_{x\mid y},\boldsymbol{\Sigma}_{x\mid y}$ & 条件分布的均值 / 协方差 & \cref{eq:cond-mean,eq:cond-cov} \\
$\mathbf{K}=\boldsymbol{\Sigma}_{xy}\boldsymbol{\Sigma}_{yy}^{-1}$ & 卡尔曼增益（雏形） & \cref{thm:gauss-cond} \\
$H(p),h(p)$ & 香农熵 / 微分熵 & \cref{def:entropy} \\
$\mathrm{KL}(p\|q)$ & KL 散度（相对熵） & \cref{def:kl} \\
$I(\mathbf{x};\mathbf{y})$ & 互信息（信息增益） & \cref{def:mutual-info} \\
\bottomrule
\end{tabular}
\end{table}

\begin{table}[H]\centering\small
\caption{定理速查表}
\begin{tabular}{lll}
\toprule
定理 / 公式 & 一句话说明 & 对应 \\
\midrule
期望线性性 \cref{eq:linearity} & 无需独立，最稳工具 & \cref{thm:linearity} \\
协方差传播 \cref{eq:affine-prop} & $\mathbf{A}\boldsymbol{\Sigma}\mathbf{A}^\top$，不确定度搬运引擎 & \cref{thm:affine-prop} \\
协方差半正定 \cref{prop:cov-psd} & $\boldsymbol{\Sigma}\succeq0$，定义高斯的前提 & \cref{prop:cov-psd} \\
全概率 \cref{eq:total-prob} & 预测步（Chapman--Kolmogorov）母公式 & \cref{def:total-prob} \\
贝叶斯定理 \cref{eq:prob-bayes} & 后验 $\propto$ 似然 $\times$ 先验 & \cref{thm:bayes} \\
高斯封闭 \cref{thm:gauss-closed} & 线性变换 / 边缘化后仍高斯 & \cref{thm:gauss-closed} \\
不相关 $\Leftrightarrow$ 独立（高斯） & 仅联合高斯下成立 & \cref{cor:gauss-uncorr-indep} \\
高斯条件分布 \cref{eq:cond-mean,eq:cond-cov} & 卡尔曼增益与更新的来源 & \cref{thm:gauss-cond} \\
高斯之积 \cref{eq:gauss-product} & 信息相加 $=$ 贝叶斯融合 & \cref{thm:gauss-product} \\
大数定律 \cref{thm:lln} & 样本均值 $\to$ 期望，蒙特卡洛基础 & \cref{thm:lln} \\
中心极限 \cref{eq:clt} & 独立小量之和 $\to$ 高斯 & \cref{thm:clt} \\
KL 非负 \cref{thm:kl-nonneg} & Gibbs 不等式，变分推断地基 & \cref{thm:kl-nonneg} \\
高斯熵 / KL \cref{eq:prob-gauss-entropy,eq:prob-gauss-kl} & $\propto\log\det\boldsymbol{\Sigma}$；信息增益尺子 & \cref{thm:gauss-info} \\
\bottomrule
\end{tabular}
\end{table}

\begin{table}[H]\centering\small
\caption{知识点总表}
\begin{tabular}{clll}
\toprule
\# & 知识点 & 核心要点 & 难度 \\
\midrule
1 & 概率空间 / $\sigma$-代数 & $(\Omega,\mathcal{F},\mathbb{P})$，三公理，连续量需可测 & \(\bigstar\bigstar\) \\
2 & 随机变量与分布 & CDF / 密度 / 质量；联合-边缘-条件 & \(\bigstar\bigstar\) \\
3 & 期望 / 方差 / 协方差 & 线性性；$\mathbf{A}\boldsymbol{\Sigma}\mathbf{A}^\top$ 传播 & \(\bigstar\bigstar\) \\
4 & 独立 vs 不相关 & 独立强于不相关；高斯下合流 & \(\bigstar\bigstar\bigstar\) \\
5 & 贝叶斯定理 & 后验 $\propto$ 似然 $\times$ 先验；基率谬误 & \(\bigstar\bigstar\bigstar\) \\
6 & 高斯分布 & 马氏距离 / 协方差椭球 / 封闭性 & \(\bigstar\bigstar\bigstar\) \\
7 & 高斯条件分布 & 舒尔补；卡尔曼增益由来 & \(\bigstar\bigstar\bigstar\bigstar\) \\
8 & 大数律 / 中心极限 & 蒙特卡洛与“高斯无处不在”的根据 & \(\bigstar\bigstar\bigstar\) \\
9 & KL / 熵 / 互信息 & 不对称尺子；信息增益 $\propto\log\det\boldsymbol{\Sigma}$ & \(\bigstar\bigstar\bigstar\bigstar\) \\
\bottomrule
\end{tabular}
\end{table}

% ════════════════════════════════════════════════════════════════
\section*{延伸阅读}
\begin{itemize}
  \item \textbf{机器人学概率范式}：Thrun、Burgard、Fox《Probabilistic Robotics》\cite{thrun2005probabilistic}——把“一切皆后验”确立为机器人学统一语言的奠基教材，贝叶斯滤波、占据栅格、定位/SLAM 的概率视角源头。
  \item \textbf{严谨主线与高斯工具}：Barfoot《State Estimation for Robotics》\cite{barfoot2024state} 第 2--3 章——概率论基元、高斯/马氏、线性高斯估计的最严谨推导（本章高斯条件分布与信息形式的完整出处）。
  \item \textbf{机器学习视角}：Bishop《Pattern Recognition and Machine Learning》\cite{bishop2006prml} 第 1--2 章——概率论、高斯分布（含条件/边缘公式的优雅证明）、信息论（熵 / KL / 互信息）的标准参考。
  \item \textbf{贝叶斯滤波与平滑}：Särkkä《Bayesian Filtering and Smoothing》\cite{sarkka2013bayesian}——从贝叶斯递推到卡尔曼/粒子滤波的系统展开，本章预测—更新框架的延伸读物。
  \item \textbf{进一步}：测度论基础（$\sigma$-代数、勒贝格积分）可参阅标准实分析/概率论教材；信息论的系统处理见 Cover \& Thomas《Elements of Information Theory》。
\end{itemize}

\section*{本章与后续章节的关系}
\begin{table}[H]\centering\small
\begin{tabular}{lll}
\toprule
后续章节 & 关系 & 本章铺垫 \\
\midrule
\cref{ch:slam_est} 状态估计 & 后验 $p(\mathbf{x}\mid\mathbf{z})$ 的 MAP / 最小二乘 & 贝叶斯、高斯、信息形式、马氏二次型 \\
\cref{ch:eskf} 卡尔曼 / ESKF & 预测—更新两步递归 & 全概率（预测）、高斯条件分布（更新、增益）\\
\cref{ch:lie} 李群 & 李群上的高斯与协方差传播 & 协方差 $\mathbf{A}\boldsymbol{\Sigma}\mathbf{A}^\top$、多元高斯 \\
\cref{ch:linalg} 线性代数 & SPD 特征分解 / 舒尔补 / 分块求逆 & 本章高斯条件分布所依赖的代数 \\
\cref{ch:convex} 凸优化 & 最大似然 / MAP 作为优化问题 & 取负对数 $=$ 二次型，凸目标 \\
\bottomrule
\end{tabular}
\end{table}

\section*{故障排查手册}
\begin{enumerate}
  \item \textbf{症状}：协方差矩阵传播后\emph{非对称或非正定}（Cholesky 崩溃）。\textbf{排查}：确认用 $\mathbf{A}\boldsymbol{\Sigma}\mathbf{A}^\top$（\cref{eq:affine-prop}）而非漏掉转置；数值上对称化 $\boldsymbol{\Sigma}\leftarrow\tfrac12(\boldsymbol{\Sigma}+\boldsymbol{\Sigma}^\top)$；条件化用 Joseph 形式保半正定（\cref{ch:eskf}）。
  \item \textbf{症状}：贝叶斯更新后后验\emph{与直觉严重不符}（如“高准确率传感器报警却几乎不更新信念”）。\textbf{原因}：忽略或设错\emph{先验}（基率谬误，\cref{ex:bayes-alarm}）。\textbf{对策}：核对先验是否反映真实发生率，单帧观测勿当确凿。
  \item \textbf{症状}：把对角协方差当“各分量独立”后，融合结果\emph{过度自信}。\textbf{原因}：非高斯下不相关 $\ne$ 独立（\cref{ex:uncorr-not-indep}）；或丢了真实的互相关 $\boldsymbol{\Sigma}_{xy}$。\textbf{对策}：保留交叉协方差块；非高斯耦合改用更细模型。
  \item \textbf{症状}：高斯条件化 / 卡尔曼增益算出\emph{异常大或 \texttt{NaN}}。\textbf{原因}：$\boldsymbol{\Sigma}_{yy}$（新息协方差）接近奇异、求逆病态。\textbf{对策}：检查观测是否退化（无信息方向）、为 $\boldsymbol{\Sigma}_{yy}$ 加观测噪声 $\mathbf{R}\succ0$ 正则。
  \item \textbf{症状}：变分 / 信息式探索里近似分布\emph{系统性偏窄或偏宽}。\textbf{原因}：KL 方向选反（\cref{def:kl} 陷阱）。\textbf{对策}：明确是 mode-seeking（$\mathrm{KL}(q\|p)$）还是 mass-covering（$\mathrm{KL}(p\|q)$），按需求选向。
  \item \textbf{症状}：把误差当高斯做外点门限，却频繁\emph{误剔正常点 / 漏检野值}。\textbf{原因}：CLT 只保中心、尾部逼近慢，真噪声可能重尾（\cref{thm:clt}）。\textbf{对策}：用卡方门限（\cref{ch:eskf}）配合鲁棒核，勿迷信 $3\sigma$。
\end{enumerate}
